\documentclass[10pt,a4paper]{article}
\usepackage[margin=2cm]{geometry}
\usepackage{fontspec}
\setmainfont{Cambria}
\setmonofont{Consolas}[Scale=0.85]
\usepackage{amsmath}
\usepackage{graphicx}
\usepackage{float}
\usepackage{longtable,booktabs,array}
\usepackage{enumitem}
\setlist{nosep,leftmargin=1.4em}
\usepackage[hidelinks]{hyperref}
\usepackage{xcolor}
\usepackage{fancyvrb}
\usepackage{newunicodechar}
\newunicodechar{∪}{\ensuremath{\cup}}
\newunicodechar{∩}{\ensuremath{\cap}}
\newunicodechar{⊗}{\ensuremath{\otimes}}
\newunicodechar{⊥}{\ensuremath{\perp}}
\newunicodechar{⊆}{\ensuremath{\subseteq}}
\newunicodechar{⊂}{\ensuremath{\subset}}
\newunicodechar{⊊}{\ensuremath{\subsetneq}}
\newunicodechar{∅}{\ensuremath{\emptyset}}
\newunicodechar{⁻}{\textsuperscript{\textminus}}
\newunicodechar{¹}{\textsuperscript{1}}
\newunicodechar{²}{\textsuperscript{2}}
\newunicodechar{⁶}{\textsuperscript{6}}
\renewcommand{\arraystretch}{1.15}
\setlength{\parskip}{3pt}
\setlength{\parindent}{0pt}
\newcolumntype{L}[1]{>{\raggedright\arraybackslash}p{#1}}

\title{CDW68 final replication report: IEEE 68-bus confirmatory replication}
\date{2026-09-13}
\author{CDW68 replication (branch research/cdw-ieee68-replication)}
\begin{document}
\maketitle
\tableofcontents
\newpage
\begin{itemize}
\item \textbf{Branch:} \texttt{research/\allowbreak{}cdw-ieee68-replication}, from \texttt{e03a5312}.
\item \textbf{Preregistration:} commit \texttt{6f188531} (2026-09-13T09:19). Every section below cites the result file it reads.
\item \textbf{Scope:} phasor-domain small-signal analysis. No EMT claim is made.
\end{itemize}
\section{Executive decision}
\textbf{CASE B}, by the exhaustive preregistered rule (prereg §17): the ranking method generalizes and the reversal does not. The gates came out as follows.

{\small
\begin{longtable}{L{0.063\linewidth}L{0.498\linewidth}L{0.379\linewidth}}
\toprule
\textbf{component} & \textbf{gate} & \textbf{result} \\
\midrule
\endhead
REV68A & level-D nested reversal in ≥ 50 \% of base-stable policies, Model A REAL & \textbf{FAIL}: 0/16 (Clopper–Pearson 95 \% upper bound 0.21) \\
REV68B & the same gate on Model B & \textbf{FAIL}: 0/16; level T also 0/16, so not PARTIAL \\
RANK68A & R12 gate, Model A & \textbf{PASS}: median ρ 0.996, advantage 0.80, top-5 1.0 \\
RANK68B & R12 gate, Model B & \textbf{PASS}: median ρ 0.981, advantage 0.55, top-5 0.8 \\
\bottomrule
\end{longtable}}
\textbf{Consequence.} The paper pivots to portfolio-conditioned dynamic reinforcement ranking. The contextual reversal is reported as IEEE-39 benchmark-specific evidence.

\textbf{The caveat that travels with the pass.} On IEEE-68 the preregistered ranking target is the rightmost transverse eigenvalue.

\begin{itemize}
\item In every portfolio of both models it is a slow real control pole: the PSS washout at −0.067 s⁻¹.
\item Branch reinforcements move it by at most 1.4e-3 s⁻¹ (A) and 3.4e-5 s⁻¹ (B), well below τ = 0.01.
\item The electromechanical version of the ranking was defined after the fact. It passes as well (A: ρ 0.998, top-5 1.0) with material effects of up to 0.047 s⁻¹, but it cannot carry the preregistered verdict.
\end{itemize}
\section{Why replication was required}
The completed hardening review left one experiment blocking a claims paper: every CDW result came from the IEEE 39-bus system with D = 0, no governors and constant-power loads. The review asked two questions:

\begin{itemize}
\item whether the nested same-mode reversal (17/19 base-stable holdout policies on IEEE-39) and the portfolio-conditioned branch ranking survive a different network with documented primary-frequency dynamics;
\item which of them survive a second converter model.
\end{itemize}
This campaign is that single test. No other network was searched.

\section{IEEE-68 provenance}
Full detail: \texttt{docs/\allowbreak{}CDW68\_\allowbreak{}MODEL\_\allowbreak{}AUDIT.md} and \texttt{results/\allowbreak{}CDW68\_\allowbreak{}MODEL\_\allowbreak{}MANIFEST.csv}.

\begin{itemize}
\item \textbf{Network} (Singh \& Pal 2013, IEEE PES TF v3.3):
\begin{itemize}
\item 68 buses, 83 branch rows, 16 of them transformers with taps; 100 MVA, 60 Hz.
\item Transcribed at Gate 3 (commit \texttt{d9fa097e}) and reproduced here: power flow within 5e-5 pu and 5e-5 deg; all 15 EM modes within 5e-4 Hz; Ybus rebuild error 0.0.
\end{itemize}
\item \textbf{Machines:} 16 sub-transient machines.
\begin{itemize}
\item Excitation: DC4B (G1–G8, G10–G12), ST1A (G9), manual (G13–G16).
\item PSS: speed washout plus three lead-lags on G1–G12.
\item Loads: constant impedance. Slack: G16.
\end{itemize}
\item \textbf{Primary frequency} (not part of the published benchmark):
\begin{itemize}
\item PST \texttt{data16m.m} (GridSTAGE copy, sha256 \texttt{c16e78f4…}) documents machine ratings (300–1900 MVA; equivalents 10–12 GVA), rotor damping d\_o (zero on G1–G12, equal to H on G13–G16) and a PST tg model-1 governor block that is commented out in the file.
\item PSTess \texttt{tg.m} (sha256 \texttt{c1584021…}) gives the equations.
\item Frozen in \texttt{inputs/\allowbreak{}pst\_\allowbreak{}primary\_\allowbreak{}frequency\_\allowbreak{}v1.json}.
\end{itemize}
\end{itemize}
\section{Dynamic model realism}
{\small
\begin{longtable}{L{0.089\linewidth}L{0.475\linewidth}L{0.080\linewidth}L{0.193\linewidth}L{0.103\linewidth}}
\toprule
\textbf{variant} & \textbf{governors} & \textbf{rotor damping} & \textbf{base α⊥ (s⁻¹)} & \textbf{rightmost EM mode} \\
\midrule
\endhead
REAL (confirmatory) & PST tg model 1 (1/R = 25 on rating, Ts 0.1, Tc 0.5, T3 0, T4 1.25, T5 5 s) on every surviving machine & PST d\_o & −0.0669 (PSS washout, real) & −0.367 s⁻¹ at 0.555 Hz \\
NOGOV (ablation B) & none & PST d\_o & −0.0461 (common-frequency mode, 0.002 Hz) & −0.371 at 0.518 Hz \\
SP33 (ablation C) & none & D = 0 (published) & −0.0556 & −0.118 at 0.520 Hz \\
\bottomrule
\end{longtable}}
\begin{itemize}
\item \textbf{Common to all variants:} equilibrium residual 7.3e-13; identical network voltages (difference 1.1e-16); smallest documented-limit margin 3.00.
\item \textbf{Damping effect.} The equivalents' documented damping moves the 0.52 Hz inter-area mode from −0.118 to −0.37 s⁻¹.
\item \textbf{Consequence for the tests.} With it, the rightmost transverse eigenvalue is a slow real pole in every portfolio. This was established in R1 before any converter portfolio was evaluated, and it is the reason the EM-tracked level T was preregistered next to level D.
\end{itemize}
\section{Converter models}
\begin{itemize}
\item \textbf{Model A: the frozen TX4 GFL} (\texttt{69f200df}; device sha256 \texttt{bcf0c76b…}).
\begin{itemize}
\item SRF-PLL 53/1400; power filter 0.03 s; outer PI 0.2/8; inner PI 0.25/6 behind x\_f 0.15, r\_f 0.01.
\item Leaky Q/V regulator with gain g (leak 0.05 rad/s).
\item Rating |S\_gen|/0.8; matched dispatch.
\item No current limits, DC link or delays.
\end{itemize}
\item \textbf{Model B: the WECC library chain TX3-GFL-0.1} (sha256 \texttt{f07a6a40…}).
\begin{itemize}
\item Blocks: PLL2, BusFreq, REGCP1, REECB1, REPCA1; constant Q; forced-PQ equilibrium.
\item Runs in ANDES 2.0.0 on a new transcription SG68D/S/M of the Singh \& Pal machine with the PST governor and damping (\texttt{code/\allowbreak{}andes68/\allowbreak{}cdw68\_\allowbreak{}models.py}).
\item Not retuned.
\end{itemize}
\item \textbf{Model B qualification} (\texttt{results/\allowbreak{}b68/\allowbreak{}B68\_\allowbreak{}qualification*.json}), all criteria met before any B portfolio was run:
\end{itemize}
{\small
\begin{longtable}{L{0.060\linewidth}L{0.894\linewidth}}
\toprule
\textbf{check} & \textbf{result} \\
\midrule
\endhead
Q1 & max |Δα⊥| against Model A's all-SG base: 9.9e-11 s⁻¹ over 16 k values, same status; EM top 3.1e-7 \\
Q2 & initialization residual ≤ 1.7e-13 \\
Q3 & descriptor QZ against reduced spectrum 7.7e-10 relative \\
Q4 & no active limiter \\
\bottomrule
\end{longtable}}
\begin{itemize}
\item \textbf{Grid-forming model:} none is validated in the repository, so none was run.
\end{itemize}
\section{Preregistration}
Files: \texttt{docs/\allowbreak{}CDW68\_\allowbreak{}PREREG\_\allowbreak{}V1.md}, \texttt{docs/\allowbreak{}CDW68\_\allowbreak{}STATISTICAL\_\allowbreak{}PLAN.md} and \texttt{docs/\allowbreak{}CDW68\_\allowbreak{}CLAIM\_\allowbreak{}MATRIX\_\allowbreak{}V1.csv}, committed at \texttt{6f188531} before any converter portfolio.

\textbf{Frozen:}

\begin{itemize}
\item the candidate set, the 16 policies, the draws and the TDS disturbance;
\item τ = 0.01 s⁻¹;
\item the level definitions A–D and T;
\item the R11/R12 gates and the exhaustive CASE rule.
\end{itemize}
\textbf{Deviations} (\texttt{docs/\allowbreak{}CDW68\_\allowbreak{}DEVIATIONS.md}, 15 entries):

\begin{itemize}
\item execution fixes (Model B check, JSON, ANDES shards, runtime registration, an R15 empty-store guard);
\item two post-hoc analyses (EM-tracked branch ranking; exploratory TDS);
\item a correction of an interpretation note (entry 11: NOGOV has three global marginals ≥ τ);
\item layout-only figure fixes;
\item paper and report production after the CASE decision (entries 13–14) and a runtime note (entry 15).
\end{itemize}
No gate, threshold, candidate, policy or model parameter was changed.

\section{Candidate selection}
\textbf{V68 = \{G3, G4, G6, G9, G11, G12\}.} These are the six physical plants with the largest scheduled active power (6.32–13.50 pu), with no tie. The rule extends the Gate 3 rule to both areas and uses no dynamic quantity.

\begin{itemize}
\item The area equivalents G13–G16 are excluded.
\item Size: 2⁶ = 64 portfolios.
\item Replacement: full, with matched dispatch.
\end{itemize}
The set was not changed after results.

\section{Policy design}
\begin{itemize}
\item \textbf{Model A: 16 maximin-LHS policies, P68\_01–16} (seed 20260913, best of 2000 designs, minimum distance 0.1748).
\begin{itemize}
\item Window: u ∈ [0, 1] with g = u², and k ∈ [0.5, 2.0].
\item P68\_01–04 are discovery and are used only to select the static comparator. P68\_05–16 are holdout.
\end{itemize}
\item \textbf{Model B: 16 conditions B68\_01–16}, carrying the k of the matching A policy.
\item All 32 conditions have a stable all-SG base, so none was excluded. No policy search took place.
\end{itemize}
\section{Portfolio census}
Output: \texttt{results/\allowbreak{}CDW68\_\allowbreak{}PORTFOLIOS.parquet} (A and B).

{\small
\begin{longtable}{L{0.206\linewidth}L{0.129\linewidth}L{0.129\linewidth}L{0.258\linewidth}L{0.219\linewidth}}
\toprule
\textbf{set} & \textbf{portfolios} & \textbf{status} & \textbf{α⊥ range (s⁻¹)} & \textbf{α⊥ in the EM band} \\
\midrule
\endhead
A REAL, 16 × 64 & 1024 & all STABLE & −0.0671 to −0.0581 & 0 \% \\
A NOGOV, 6 × 64 & 384 & all STABLE & −0.0612 to −0.0461 & 0 \% \\
A SP33, 6 × 64 & 384 & all STABLE & −0.0622 to −0.0539 & 0 \% \\
B REAL, 16 × 64 & 1024 & all STABLE & −0.06695 to −0.06688 & 0 \% \\
A draws, 40 × 64 & 2560 & all STABLE & max −0.0669 & 0 \% \\
B draws, 20 × 64 & 1280 & all STABLE & max −0.0669 & 0 \% \\
\bottomrule
\end{longtable}}
\begin{itemize}
\item No portfolio is unstable, so H\_RHP is empty and κ is undefined in every condition.
\item \textbf{Rightmost EM-band mode:}
\begin{itemize}
\item A REAL: 0.550–0.806 Hz (median 0.557 Hz), real part −0.373 to −0.263 s⁻¹;
\item B: 0.552–0.635 Hz, real part −0.369 to −0.251 s⁻¹.
\end{itemize}
\end{itemize}
\section{Contextual reversal}
\textbf{Primary result} (\texttt{results/\allowbreak{}CDW68\_\allowbreak{}R07\_\allowbreak{}A\_\allowbreak{}gate.json}, \texttt{results/\allowbreak{}CDW68\_\allowbreak{}R07\_\allowbreak{}B\_\allowbreak{}gate.json}).

\begin{itemize}
\item Level-D nested reversal coverage is \textbf{0/16 for Model A} and \textbf{0/16 for Model B} (95 \% Clopper–Pearson upper bound 0.21 each).
\item Levels A, B and C are also 0/16, at every τ from 0.01 to 0.05.
\item The median-magnitude clause is undefined, because no level-D pair exists.
\end{itemize}
\textbf{Why levels A–D are empty.} The global marginals move a slow real pole.

\begin{itemize}
\item The largest level-A |Δα⊥| is 8.7e-3 s⁻¹ (A REAL, P68\_16, replacing G12). For B it is 2.3e-5.
\item No global marginal is material, so no sign class is non-zero.
\end{itemize}
\textbf{Preregistered secondary, level T (tracked rightmost EM mode):} 0/16 for both models.

{\small
\begin{longtable}{L{0.060\linewidth}L{0.171\linewidth}L{0.154\linewidth}L{0.060\linewidth}L{0.248\linewidth}L{0.256\linewidth}}
\toprule
\textbf{model} & \textbf{EM-tracked marginals} & \textbf{range (s⁻¹)} & \textbf{median} & \textbf{materially stabilizing (≤ −τ)} & \textbf{materially destabilizing (≥ τ)} \\
\midrule
\endhead
A REAL & 2770 & −0.0069 to +0.0528 & +0.0020 & 0 & 437 \\
B REAL & 2804 & −0.0008 to +0.0570 & +0.0101 & 0 & 1429 \\
\bottomrule
\end{longtable}}
Replacing a synchronous machine by either converter never materially stabilizes the tracked electromechanical mode on this network, in any context. A reversal needs a stabilizing side of at least τ, so none can exist. F2 and F3 show the one-sided pattern.

\section{Discrete curvature}
\begin{itemize}
\item \textbf{R8.} Chain-identity verification is vacuous: no nested reversal exists at any level.
\item \textbf{Interaction terms still exist} (F4). On the rightmost EM mode, second differences have both signs beyond 1e-3 s⁻¹ in 13/16 A REAL policies and in 16/16 B policies. Their range is −0.0067 to +0.0457 s⁻¹ (A) and −0.0045 to +0.0562 s⁻¹ (B).
\item \textbf{Reference policy P68\_10:} 240 one-step squares with range −0.0052 to +0.0009 s⁻¹; 104 are below −1e-3 and none is above +1e-3. On the global α⊥ all |d| ≤ 2.9e-5.
\item \textbf{Reading.} The set function is not modular, but its interactions do not change the sign of a replacement's effect: marginals shrink or grow with context and stay destabilizing.
\item \textbf{Post-hoc TDS pairs} (§18). On their single-step chains the identity holds exactly (residual 0).
\end{itemize}
\section{Fixed ranking}
Source: \texttt{results/\allowbreak{}CDW68\_\allowbreak{}R09\_\allowbreak{}gate.json}.

{\small
\begin{longtable}{L{0.060\linewidth}L{0.537\linewidth}L{0.356\linewidth}}
\toprule
\textbf{stratum} & \textbf{Model A} & \textbf{Model B} \\
\midrule
\endhead
FULL & 57 decisions per policy (median), 0 material preferences: p* undefined, regret 0 & same \\
EM & 0 decisions: not evaluable & same \\
LEVEL-D & 0 decisions: not evaluable & same \\
TRACKED & 51.5 decisions, p* = 1.0, regret 0, transfer regret 0 & 46.5 decisions, p* = 1.0, regret 0, transfer regret 0 \\
\bottomrule
\end{longtable}}
\begin{itemize}
\item The insufficiency bar (p* ≤ 0.90 and regret ≥ 0.10) is not met in any stratum.
\item The T1 sign test cannot be computed.
\item The headline "fixed ranking insufficient" is not allowed.
\item On IEEE-68 the oracle fixed order of units reproduces every material preference on the tracked EM mode.
\end{itemize}
\section{Total sensitivity validation (R11)}
Source: \texttt{results/\allowbreak{}CDW68\_\allowbreak{}R12\_\allowbreak{}gate.json}.

{\small
\begin{longtable}{L{0.060\linewidth}L{0.381\linewidth}L{0.060\linewidth}L{0.067\linewidth}L{0.101\linewidth}L{0.060\linewidth}L{0.313\linewidth}}
\toprule
\textbf{model} & \textbf{test} & \textbf{n} & \textbf{sign agreement} & \textbf{median relative error} & \textbf{Spearman} & \textbf{verdict} \\
\midrule
\endhead
A & IFT total derivative (SPR-68) vs central re-equilibrated difference at h = 1e-4 & 80 & 1.000 & 1.3e-5 & 0.99995 & \textbf{PASS} \\
B & numerical total derivative, h = 1e-3 vs h = 1e-4 & 80 & 1.000 & 5.1e-7 & 0.99991 & \textbf{PASS} (numerical consistency only; B has no analytical engine) \\
\bottomrule
\end{longtable}}
\section{Reinforcement ranking (R12)}
Holdout: 12 conditions per model, γ = 1.5.

{\small
\begin{longtable}{L{0.396\linewidth}L{0.307\linewidth}L{0.237\linewidth}}
\toprule
\textbf{} & \textbf{Model A} & \textbf{Model B} \\
\midrule
\endhead
median Spearman(D\_tot, finite) & 0.9955 & 0.9810 \\
selected static (discovery) & S4 electrical distance, ρ 0.197 & S9 ΔgSCR, ρ 0.434 \\
median paired advantage (95 \% bootstrap) & 0.799 (0.797–0.800) & 0.547 (0.546–0.547) \\
median top-5 precision & 1.0 & 0.8 \\
median Kendall / NDCG@5 & 0.958 / 1.000 & 0.896 / 0.998 \\
ρ at γ = 1.10 / 1.25 / 1.50 & 0.9993 / 0.9980 / 0.9955 & 0.9984 / 0.9939 / 0.9810 \\
T2 sign-flip p (Holm) & 0.0078 (0.0156) & 0.0898 (0.0898) \\
Wilcoxon p & 0.00024 & 0.00024 \\
gate & \textbf{PASS} & \textbf{PASS} \\
\bottomrule
\end{longtable}}
\begin{itemize}
\item \textbf{T2B.} It does not reach 0.05. The preregistered statistic is the median under exact sign-flip enumeration. The twelve B advantages are nearly equal in size, and in that case the median statistic has little power. The gate itself is defined on medians and passes.
\item \textbf{Materiality.} The median over conditions of the largest finite effect is 4.2e-5 s⁻¹ (A) and 3.3e-5 s⁻¹ (B). In the median condition no branch has a material effect.
\item \textbf{Post-hoc EM-tracked ranking (A only).} Median ρ 0.998 and top-5 1.0. The median of the per-condition largest effect is 0.020 s⁻¹ (maximum 0.047). A median of 5 branches per condition exceed τ.
\item \textbf{Top branches.}
\begin{itemize}
\item EM-tracked, by median finite effect: lines 18–50, 51–50 and 18–49.
\item Preregistered target (both models): the G1 step-up transformer 1–54.
\end{itemize}
\end{itemize}
\section{Static baselines}
Discovery medians of Spearman with the finite ×1.5 effect (\texttt{results/\allowbreak{}CDW68\_\allowbreak{}R10\_\allowbreak{}static.csv}):

{\small
\begin{longtable}{L{0.527\linewidth}L{0.206\linewidth}L{0.206\linewidth}}
\toprule
\textbf{index} & \textbf{A} & \textbf{B} \\
\midrule
\endhead
S1 |P| & 0.043 & 0.202 \\
S2 |S| & 0.047 & 0.221 \\
S3 |z| & 0.021 & 0.230 \\
S4 electrical distance & \textbf{0.194} & 0.337 \\
S5 effective resistance & 0.116 & 0.294 \\
S6 Fiedler edge & 0.093 & 0.242 \\
S7 betweenness & 0.029 & 0.081 \\
S8 dV/dQ & 0.123 & −0.014 \\
S9 ΔgSCR & 0.067 & \textbf{0.434} \\
\bottomrule
\end{longtable}}
\begin{itemize}
\item The static indices do not change across conditions.
\item In the holdout, the best static index per condition has median ρ 0.197 (A) and 0.434 (B).
\item The Laplacian scores (S5, S6) are baselines only; no topology explanation is offered.
\end{itemize}
\section{Governor/damping ablation (R13)}
Model A, P68\_01–06, 64 portfolios each.

{\small
\begin{longtable}{L{0.060\linewidth}L{0.060\linewidth}L{0.060\linewidth}L{0.153\linewidth}L{0.104\linewidth}L{0.118\linewidth}L{0.418\linewidth}}
\toprule
\textbf{variant} & \textbf{level D} & \textbf{level T} & \textbf{EM-tracked range (s⁻¹)} & \textbf{stabilizing ≥ τ} & \textbf{destabilizing ≥ τ} & \textbf{largest global |Δα⊥|} \\
\midrule
\endhead
REAL & 0/6 & 0/6 & −0.0052 to +0.0393 & 0 & 267 & 1.5e-4 \\
NOGOV & 0/6 & 0/6 & −0.0102 to +0.0228 & 2 & 271 & 1.10e-2 (3 stabilizing, common-frequency mode, level B only) \\
SP33 & 0/6 & 0/6 & −0.0136 to +0.0274 & 6 & 277 & 4.7e-3 \\
\bottomrule
\end{longtable}}
\begin{itemize}
\item Removing governors, or both governors and damping, makes a few EM-tracked marginals materially stabilizing. None of them forms a nested reversal.
\item Reversal is absent in the realistic model and in both simplifications. On this network the documented primary-frequency dynamics neither create nor remove it.
\item The simplification does shift the EM marginals slightly toward stabilizing values, which is the direction the IEEE-39 limitation paragraph anticipated.
\end{itemize}
\section{Cross-converter transfer}
\begin{itemize}
\item \textbf{What transfers between A and B:}
\begin{itemize}
\item the absence of reversal (0/16 in both);
\item one-sided destabilizing EM effects (B's are larger: median +0.0101 against +0.0020);
\item an error-free fixed ranking on the tracked mode;
\item the ranking method (R12 PASS in both).
\end{itemize}
\item \textbf{What does not transfer: the line list.} Across matched conditions, the Spearman between the two models' finite truths has median 0.50 (range 0.00–0.51), Kendall 0.37 and top-5 overlap 0.4. The 0.00 is P68\_16, where Model A's critical pole is a different pole.
\item \textbf{Which static index serves as comparator also depends on the converter:} electrical distance for A, ΔgSCR for B.
\end{itemize}
\section{Nonlinear TDS}
Source: \texttt{results/\allowbreak{}CDW68\_\allowbreak{}R14\_\allowbreak{}tds\_\allowbreak{}posthoc.json}, figure F10.

\textbf{As preregistered.} No case exists: there is no level-D or level-T reversal and no level-D-eligible negative control. R14 is "not applicable".

\textbf{Post-hoc exploratory run.} The run took the two EM-tracked nested pairs with the largest same-sign marginals at the reference policy P68\_10.

\begin{itemize}
\item Model A REAL; 0.5 pu reactor at bus 3 for 10 s; decay of the tracked modal coordinate fitted over 11–29 s.
\item Integrator: BDF with a network solve at each step.
\end{itemize}
{\small
\begin{longtable}{L{0.209\linewidth}L{0.274\linewidth}L{0.274\linewidth}L{0.183\linewidth}}
\toprule
\textbf{pair} & \textbf{TDS decay change} & \textbf{linear decay change} & \textbf{sign agreement} \\
\midrule
\endhead
G12, BASE → \{11\} & +0.011463 / +0.011467 & +0.011460 / +0.011466 & 2/2 \\
G4, \{3\} → \{3,9\} & +0.010275 / +0.010310 & +0.010273 / +0.010307 & 2/2 \\
\bottomrule
\end{longtable}}
The traces confirm the linear signs and magnitudes to within 4e-6 s⁻¹. They say nothing about reversals, because none exists.

\section{Uncertainty}
Deterministic stress envelopes E05 (±5 \%) and E10 (±10 \%) at the reference policy (\texttt{results/\allowbreak{}CDW68\_\allowbreak{}R15\_\allowbreak{}uncertainty.json}).

{\small
\begin{longtable}{L{0.060\linewidth}L{0.060\linewidth}L{0.105\linewidth}L{0.105\linewidth}L{0.322\linewidth}L{0.158\linewidth}L{0.184\linewidth}}
\toprule
\textbf{set} & \textbf{draws} & \textbf{level-D coverage} & \textbf{level-T coverage} & \textbf{draws with any materially stabilizing EM marginal} & \textbf{all 64 portfolios stable} & \textbf{smallest EM-tracked marginal} \\
\midrule
\endhead
A E05 & 20 & 0/20 & 0/20 & 0/20 & 20/20 & −0.0055 \\
A E10 & 20 & 0/20 & 0/20 & 0/20 & 20/20 & −0.0056 \\
B E05 & 10 & 0/10 & 0/10 & 0/10 & 10/10 & −0.0009 \\
B E10 & 10 & 0/10 & 0/10 & 0/10 & 10/10 & −0.0014 \\
\bottomrule
\end{longtable}}
\begin{itemize}
\item \textbf{Witness identity} is undefined: the nominal reference has no witness and no draw creates one.
\item \textbf{Ranking under draws} (Model A; first 5 draws of each envelope, 83 branches, ×1.5): median ρ of D\_tot 0.996 (10–90 \% range over draws 0.9955–0.9960); post-hoc EM-tracked ρ 0.998 (0.9974–0.9981); median largest finite effect 4.3e-5 s⁻¹.
\end{itemize}
These are fractions of designed draws, not probabilities.

\section{Cross-benchmark comparison}
Source: \texttt{results/\allowbreak{}CDW\_\allowbreak{}CROSS\_\allowbreak{}BENCHMARK\_\allowbreak{}MATRIX.csv}, figure F12.

{\small
\begin{longtable}{L{0.216\linewidth}L{0.208\linewidth}L{0.216\linewidth}L{0.158\linewidth}L{0.141\linewidth}}
\toprule
\textbf{} & \textbf{IEEE-39 A} & \textbf{IEEE-39 B} & \textbf{IEEE-68 A} & \textbf{IEEE-68 B} \\
\midrule
\endhead
governors / damping & none / D = 0 & none / D = 0 & PST tg / PST d\_o & PST tg / PST d\_o \\
candidates, policies & 9; 19 base-stable holdout & 4; 7 & 6; 16 & 6; 16 \\
level-D reversal & 17/19 (median 0.0127 s⁻¹) & global 0/7; EM-tracked 1/7 & 0/16 (T 0/16) & 0/16 (T 0/16) \\
fixed-ranking regret & FULL 0.38; EM 0.03 & not computed & FULL 0; TRACKED 0 & FULL 0; TRACKED 0 \\
total-sensitivity median ρ & 0.996 & 0.993 & 0.996 & 0.981 \\
static comparator median ρ & 0.445 & 0.790 & 0.197 & 0.434 \\
advantage & 0.49 & 0.20 & 0.80 & 0.55 \\
top-5 & 1.0 & n/a & 1.0 & 0.8 \\
TDS & not performed & not performed & post hoc, signs 4/4 & not applicable \\
\bottomrule
\end{longtable}}
\section{Negative results}
All are preserved in the result files.

\begin{enumerate}
\item The primary reversal gate fails on both models: 0/16 and 0/16.
\item Level T fails on both models: 0/16 and 0/16.
\item No materially stabilizing EM-tracked marginal exists in REAL for either model, in the census or in any of the 60 draws.
\item Ablations: 0/6 reversals without governors and 0/6 without governors or damping.
\item The fixed-ranking failure of IEEE-39 does not replicate. Every IEEE-68 stratum has regret 0 or cannot be evaluated.
\item The preregistered branch-ranking target is immaterial: its effects are ≤ 1.4e-3 s⁻¹, and in the median condition no branch is material.
\item T2B is not significant (sign-flip p 0.090).
\item The line list does not transfer between converter models (median ρ 0.50).
\item The preregistered TDS holdout has no case.
\item Entry 7 of the deviation log was wrong for NOGOV and is corrected in entry 11.
\end{enumerate}
\section{Novelty implications}
See \texttt{docs/\allowbreak{}CDW68\_\allowbreak{}NOVELTY\_\allowbreak{}ASSESSMENT.md}.

\begin{itemize}
\item Nested contextual reversal and fixed-ranking failure become \textbf{IEEE-39 benchmark-specific evidence}.
\item The one finding carried across two benchmarks and two converter models is the portfolio-conditioned total branch sensitivity used as a ranking. It is a known construction (Smed 1993; Nam et al. 2000; Joswig-Jones et al. 2026, preprint), validated against finite reinforcements on held-out policies and compared with discovery-selected static indices.
\item On IEEE-68 its preregistered target is immaterial. The material EM version is post hoc.
\item The word "first" is not used.
\end{itemize}
\section{Paper decision}
\textbf{CASE B.} The CDW manuscript (\texttt{reports/\allowbreak{}papers/\allowbreak{}cdw\_\allowbreak{}contextual\_\allowbreak{}dynamic\_\allowbreak{}weakness/\allowbreak{}}) is re-titled and its contributions reordered:

\begin{enumerate}
\item Portfolio-conditioned re-equilibrated reinforcement ranking on two benchmarks and two converter models.
\item When re-equilibration matters: on IEEE-68 the frozen derivative reaches ρ 0.55 against 0.996; total and frozen disagree in sign for a median 33 \% of branches.
\item Contextual reversal on IEEE-39, with its failure to replicate on IEEE-68.
\item Boundaries: fixed rankings, cross-model transfer, and the immaterial IEEE-68 target.
\end{enumerate}
\textbf{Readiness.} The manuscript is not yet ready for TPWRS as a claims paper. The IEEE-68 ranking pass rests on an immaterial pole. Its material EM version is post hoc and needs a preregistered confirmation, or at least a second EM-target test on a fresh set of conditions.

\section{Reproducibility}
\begin{itemize}
\item \textbf{Software:} Python 3.13; NumPy, SciPy and pandas from \texttt{.venv/\allowbreak{}tx3-analysis}; ANDES 2.0.0 in \texttt{.venv/\allowbreak{}xtool-andes-gfl}. Thread counts are fixed to 1.
\item \textbf{Code} in \texttt{code/\allowbreak{}}:
\end{itemize}
{\small
\begin{longtable}{L{0.254\linewidth}L{0.686\linewidth}}
\toprule
\textbf{phase} & \textbf{scripts} \\
\midrule
\endhead
inputs and audit & \texttt{R00\_\allowbreak{}inputs.py}, \texttt{R01\_\allowbreak{}audit.py}, \texttt{R01b\_\allowbreak{}manifest.py} \\
Model B & \texttt{R02\_\allowbreak{}export\_\allowbreak{}b.py}, \texttt{andes68/\allowbreak{}run68.py} \\
census and reversal & \texttt{R06\_\allowbreak{}census.py}, \texttt{R07\_\allowbreak{}reversal.py} \\
fixed ranking & \texttt{R09\_\allowbreak{}ranking.py} \\
branches and gates & \texttt{R10\_\allowbreak{}branches.py}, \texttt{R10\_\allowbreak{}static.py}, \texttt{R12\_\allowbreak{}ranking\_\allowbreak{}gate.py} \\
TDS, uncertainty, synthesis & \texttt{R14\_\allowbreak{}tds.py}, \texttt{R15\_\allowbreak{}uncertainty.py}, \texttt{R18\_\allowbreak{}crossbench.py}, \texttt{R21\_\allowbreak{}figures.py} \\
\bottomrule
\end{longtable}}
\begin{itemize}
\item \textbf{Raw per-task JSON} in \texttt{raw/\allowbreak{}} is not committed. It is listed with sha256 in \texttt{CDW68\_\allowbreak{}RAW\_\allowbreak{}MANIFEST.csv} and packed in \texttt{CDW68\_\allowbreak{}RAW\_\allowbreak{}RESULTS\_\allowbreak{}20260913.zip}.
\item \textbf{Runtimes (wall clock):}
\end{itemize}
{\small
\begin{longtable}{L{0.593\linewidth}L{0.347\linewidth}}
\toprule
\textbf{phase} & \textbf{time} \\
\midrule
\endhead
census, 1792 cases & 451 s \\
Model B census & 8 shards, 55 s \\
Model B branches & 6 shards, ≤ 919 s \\
Model B draws & 5 shards, ≤ 160 s \\
Model A branches, 128 tasks & 1862 s \\
Model A draws, 2560 cases & 703 s \\
ranking under draws, 40 tasks & 712 s \\
\bottomrule
\end{longtable}}
\begin{itemize}
\item The campaign ran from the preregistration commit (09:19) through the last computation (R15L, finished 10:48:02) on 2026-09-13; the times of the document, paper and bundle commits are in git.
\item \textbf{Figures:} F1–F12 in \texttt{figures/\allowbreak{}} (PDF, SVG, PNG).
\end{itemize}
\end{document}